% !TEX program = pdflatex
% CARNET FICHE 7 — THEME 2 — Equation d'etat pV=nRT — TUTO
\documentclass[11pt,a4paper]{article}
\usepackage[utf8]{inputenc}
\usepackage[T1]{fontenc}
\usepackage{lmodern}
\usepackage[french]{babel}
\usepackage{amsmath,amssymb,mathtools}
\usepackage{icomma}
\usepackage{siunitx}
\sisetup{output-decimal-marker={,},per-mode=symbol}
\DeclareSIUnit{\atm}{atm}
\DeclareSIUnit{\bar}{bar}
\DeclareSIUnit{\mmHg}{mmHg}
\usepackage[version=4]{mhchem}
\usepackage{xcolor}
\usepackage{colortbl}
\usepackage{tikz}
\usepackage[most]{tcolorbox}
\usepackage{enumitem}
\usepackage{array}
\usepackage{booktabs}
\usepackage{fancyhdr}
\usepackage[a4paper,margin=2.1cm,headheight=26pt,headsep=10pt]{geometry}
\usepackage{titlesec}
\usetikzlibrary{arrows.meta,positioning,calc,shapes.geometric,fit}

\definecolor{primary}{RGB}{150,98,162}
\definecolor{primarydark}{RGB}{92,52,110}
\definecolor{defcol}{RGB}{0,114,178}
\definecolor{formulacol}{RGB}{198,93,9}
\definecolor{attncol}{RGB}{200,40,55}
\definecolor{methcol}{RGB}{90,90,100}
\definecolor{pcol}{RGB}{200,40,55}
\definecolor{vcol}{RGB}{0,135,90}
\definecolor{tcol}{RGB}{198,93,9}
\definecolor{ncol}{RGB}{120,94,200}
\definecolor{gazcol}{RGB}{0,120,190}

\tcbset{boxbase/.style={enhanced,breakable,boxrule=0.9pt,arc=2.5pt,
   left=3mm,right=3mm,top=2.2mm,bottom=2.2mm,fonttitle=\bfseries,coltitle=white}}
\newtcolorbox{idee}[1][]{boxbase,colback=defcol!6,colframe=defcol,title={#1}}
\newtcolorbox{formule}[1][]{boxbase,colback=formulacol!8,colframe=formulacol,title={#1}}
\newtcolorbox{piege}[1][]{boxbase,colback=attncol!7,colframe=attncol,title={#1}}
\newtcolorbox{methode}[1][]{boxbase,colback=methcol!8,colframe=methcol,sharp corners,title={#1}}

\pagestyle{fancy}\fancyhf{}
\fancyhead[L]{\small\textcolor{primarydark}{\textbf{Fiche 7} — Loi des gaz}}
\fancyhead[R]{\small\textcolor{primarydark}{\textsc{Thème 2 · Tuto}}}
\fancyfoot[C]{\small\thepage}
\renewcommand{\headrulewidth}{0.8pt}
\renewcommand{\headrule}{\hbox to\headwidth{\color{primary}\leaders\hrule height \headrulewidth\hfill}}
\titleformat{\section}{\large\bfseries\color{primarydark}}{\thesection}{0.6em}{}
\setlength{\parindent}{0pt}\setlength{\parskip}{4pt}

\begin{document}

\begin{center}
{\Large\bfseries\color{primarydark} Thème 2 — L'équation d'état $pV=nRT$}\\[2pt]
{\color{primary}\itshape Tuto à lire avant de commencer les exercices}
\end{center}
\vspace{-2pt}
\textcolor{primary}{\hrule height 1pt}
\vspace{6pt}

\section*{1. Une seule équation pour les quatre variables}

Au Thème 1, on comparait \emph{deux états} d'un même gaz par des rapports. L'équation d'état, elle, relie
les quatre variables \textbf{dans un seul état}, à un instant donné :

\begin{formule}[Loi des gaz parfaits]
\[ \boxed{\;\textcolor{pcol}{p}\,\textcolor{vcol}{V}=\textcolor{ncol}{n}\,R\,\textcolor{tcol}{T}\;} \]
\end{formule}
où \textcolor{pcol}{$p$} est la pression, \textcolor{vcol}{$V$} le volume, \textcolor{ncol}{$n$} la
quantité de matière (en moles), \textcolor{tcol}{$T$} la température \textbf{absolue (kelvins)} et $R$ la
\textbf{constante des gaz parfaits}. Connaissant trois des quatre grandeurs, on trouve la quatrième.

\section*{2. La constante $R$ : choisir la valeur cohérente avec ses unités}

$R$ est universelle (la même pour tous les gaz), mais sa \emph{valeur numérique} dépend des unités de $p$
et de $V$. \textbf{On choisit une ligne du tableau et on s'y tient} :

\begin{center}\renewcommand{\arraystretch}{1.35}\small
\begin{tabular}{c c c c}
\toprule
\rowcolor{primary!12}
\textbf{Valeur de $R$} & \textbf{$p$ en} & \textbf{$V$ en} & \textbf{$T$ en}\\
\midrule
$R=\qty{0.082}{\atm\litre\per\mole\per\kelvin}$ & \si{\atm} & \si{\litre} & \si{\kelvin}\\
$R=\qty{8.314}{\joule\per\mole\per\kelvin}$ & \si{\pascal} & \si{\meter\cubed} & \si{\kelvin}\\
$R=\qty{8.314}{\kilo\pascal\litre\per\mole\per\kelvin}$ & \si{\kilo\pascal} & \si{\litre} & \si{\kelvin}\\
\bottomrule
\end{tabular}
\end{center}

\begin{piege}[L'erreur la plus fréquente]
\textbf{Ne jamais panacher.} Une pression en \si{\atm} avec un volume en \si{\litre} $\Rightarrow$
$R=\num{0.082}$. Utiliser $R=8{,}314$ avec des atmosphères fausse le résultat d'un facteur $\sim100$.
Et \emph{toujours} $T$ en kelvins ($T=\theta+273$).
\end{piege}

\section*{3. Les quatre formes de la même équation}

Selon l'inconnue, on isole la grandeur cherchée \emph{avant} de calculer :
\begin{center}\small
$\displaystyle n=\frac{pV}{RT}$ \qquad
$\displaystyle p=\frac{nRT}{V}$ \qquad
$\displaystyle V=\frac{nRT}{p}$ \qquad
$\displaystyle T=\frac{pV}{nR}$
\end{center}
La plus utilisée en chimie est $n=\dfrac{pV}{RT}$ : elle donne le nombre de moles à partir des conditions
$(p,V,T)$, ce qui ouvre ensuite la porte aux masses et à la stœchiométrie.

\begin{idee}[Le pont vers la masse : $n=\dfrac{m}{M}$]
L'équation ne contient pas la masse. Mais comme $n=\dfrac{m}{M}$ (Fiche 4), on passe de la quantité de
matière à la \textbf{masse} $m$ (ce qu'on pèse) et inversement. C'est ce chaînon qui permet de
« peser un gaz » : $m=n\times M$ après avoir trouvé $n=\dfrac{pV}{RT}$.
\end{idee}

\section*{4. Deux exemples types}

\textbf{(a) Trouver une pression.} On enferme $n=\qty{0.50}{\mole}$ de \ce{O2} dans un récipient rigide
$V=\qty{5.0}{\litre}$ à $\theta=\qty{27}{\degreeCelsius}$. D'abord $T=\qty{300}{\kelvin}$. On isole $p$ :
\[ p=\frac{nRT}{V}=\frac{0{,}50\times0{,}082\times300}{5{,}0}=\qty{2.46}{\atm}. \]

\textbf{(b) Comparer deux enceintes à la même température.} Deux gaz sont à la même température $T$. Alors
$\dfrac{pV}{n}=RT$ vaut la \emph{même chose} pour les deux, d'où
$\dfrac{p_1V_1}{n_1}=\dfrac{p_2V_2}{n_2}$. Nul besoin de connaître $T$ ni $R$ : on raisonne en rapports,
exactement comme au Thème 1, mais entre deux gaz différents.

\begin{methode}[Démarche pour un problème avec $pV=nRT$]
\textbf{1.} Repérer l'inconnue et les trois données. \quad
\textbf{2.} Convertir : $T$ en kelvins ; choisir $R$ cohérent avec les unités de $p$ et $V$. \quad
\textbf{3.} Isoler littéralement l'inconnue. \quad
\textbf{4.} Si l'on cherche une masse : trouver $n$, puis $m=nM$. \quad
\textbf{5.} Calculer et vérifier l'ordre de grandeur (une mole de gaz $\approx$ quelques dizaines de
grammes, quelques dizaines de litres aux conditions usuelles).
\end{methode}

\end{document}
